\documentclass{article}
\usepackage{amsmath}
\usepackage{cancel}





\title{Taylor Series in Two Dimensions}
\begin{document}
Let's try and extend our notion of a Taylor series for a function $f$ of a
single variable $x$ to two dimensions. To that end, let's now consider functions of
the form
\begin{equation}
  f = f(x_1, x_2) 
\end{equation}
The properties of polynomials are \textit{really nice} and therefore it would be
nice to have some kind of polynomial approximation to our function $f$. That is,
we want to look for an approximation of the form
\begin{equation}
  f(x_1, x_2) = \sum_{i=0}^{\infty}\sum_{j=0}^{\infty} c_{ij}(x_1-a)^i(x_2-b)^j
\end{equation}
about the point $(x_1 = a, x_2 = b)$.

Our task is to find the expansion coefficients $c_{ij}$. Again, one of the nice
properties of polynomials is they are really easy to differentiate. Let's do
this a couple of times and evaluate the result at the point $(a,b)$.
\begin{align}
  &f(x_1,x_2) = c_{00} + c_{10}(x_1-a) + c_{01}(x_2-b)+ c_{11}(x_1-a)(x_2-b) + c_{20}(x_1-a)^2+ c_{02}(x_2-b)^2 + ... \\
  &\Rightarrow f(a,b) = c_{00} + 0 + 0 + 0 + ... \\
  &\frac{\partial f}{\partial x_1}(x_1, x_2) = c_{10}  + c_{11}(x_2-b) + 2c_{20}(x_1-a) + ... \\
  &\Rightarrow f_{x_1}(a,b) = c_{10} + 0 + 0 + ... \\
  &\frac{\partial f}{\partial x_2}(x_1,x_2) = c_{01} + c_{11}(x_1-a) + 2c_{02}(x_2-b) + ... \\
  &\Rightarrow f_{x_2}(a,b) = c_{01} + 0 + 0  + ... \\
  &\frac{\partial^2 f}{\partial x_1 \partial x_1} = 2c_{20}+ 2c_{22}(x_2-b)^2 + 6c_{30}(x_1-a) + ... \\
  &\quad \vdots  \hspace{5em}\vdots \quad \textrm{continuing the pattern} \\
  &\quad \vdots  \hspace{5em}\vdots \\ 
  &f_{x_1x_1}(a,b) = 2c_{20} + 0 + 0 + ... \\
  &f_{x_2x_2}(a,b) = 2c_{02} + 0 + 0 + ... \\
  &f_{x_1x_2}(a,b) = c_{11} + 0 + 0 + ... \\
  &f_{x_2x_1}(a,b) = c_{11} + 0 + 0 + ...  \\ 
  &\quad \vdots  \hspace{5em}\vdots 
\end{align}
Let's tabulate our results so far. We have
\begin{center}
 \begin{tabular}{||c c c ||} 
   \hline
   coefficient & value & k$\equiv$i+j \\ 
   \hline\hline
   $c_{00}$ & $f(a,b)$ & $0$ \\ 
   \hline
   $c_{10}$ & $f_{x_1}(a,b)$ & $1$ \\
   \hline
   $c_{01}$ & $f_{x_2}(a,b)$ & $1$ \\
   $c_{20}$ & $\frac{1}{2}f_{x_1 x_1}(a,b)$ & $2$ \\
   $c_{02}$ & $\frac{1}{2}f_{x_2 x_2}(a,b)$ & $2$ \\
   $c_{11}$ & $\frac{1}{2}\left(f_{x_1 x_2}(a,b) + f_{x_2 x_1}(a,b)\right)$ & $2$ 
\end{tabular}
\end{center}
Let's step back and look at this pattern. Remember that for the 1-D Taylor
series, we have
\begin{align}
  f(z) &= \sum_{n=0}^\infty c_n(x-a)^n \\
  c_n &= \frac{f^{(n)}(a)}{n!}
\end{align}
The pattern appears to be the same for the 2-D Taylor series for the combined
order $k \equiv i + j$. The only difference is we have many different ways to
combine $i$ $x_1$ derivatives with $j$ $x_2$ derivatives to make $k$ derivatives
in total. Furthermore, because mixed partial derivatives are equivalent
(according to Clairaut's theorem), we can take advantage of binomial
coefficients to write the general form
\begin{equation}
  c_{ij} = \frac{1}{(i+j)!}\frac{\partial^i}{\partial x_1^i}\frac{\partial^j}{\partial x_2^j}f(a,b)\;\begin{pmatrix}i+j\\ i\end{pmatrix}(x_1-a)^i(x_2-b)^j
\end{equation}




\end{document}
